\subsection{Density, Separability, and Closure Tests}

\begin{proposition} \label{prop:dense-closure-criterion} \hlt{Density and Closure}\\
Let $X$ be a metric space and let $E\subseteq X$. Then $E$ is dense in $X$ if and only if $\overline{E}=X$.
\begin{proof}
\begin{enumerate}[label=\roman*.]
\setlength{\itemsep}{0pt}
\item[$(\Rightarrow)$] Suppose $E$ is dense in $X$. Let $p\in X$. By Definition \ref{def:metric-topology-core-definitions}, $p$ is either a point of $E$ or a limit point of $E$. Hence $p\in E\cup E^{\operatorname{acc}}=\overline{E}$ by Definition \ref{def:metric-closure}. Thus $X\subseteq\overline{E}$. Since $\overline{E}\subseteq X$, $\overline{E}=X$.
\item[$(\Leftarrow)$] Suppose $\overline{E}=X$. Let $p\in X$. Then $p\in\overline{E}=E\cup E^{\operatorname{acc}}$, so $p$ is either a point of $E$ or a limit point of $E$. Therefore $E$ is dense in $X$ by Definition \ref{def:metric-topology-core-definitions}.
\end{enumerate}
\hlt{Structural Commentary}: Density says every point is already in $E$ or is forced by nearby points of $E$.\\
\hlt{Key Insight}: The closure packages exactly those two possibilities.
\end{proof}
\end{proposition}

\begin{proposition} \label{prop:dense-closure-and-interval-test} \hlt{Real-Line Density Interval Test}\\
Let $E\subseteq\R$. Then $E$ is dense in $\R$ if and only if every nonempty open interval $(a,b)$ satisfies $(a,b)\cap E\neq\emptyset$.
\begin{proof}
\begin{enumerate}[label=\roman*.]
\setlength{\itemsep}{0pt}
\item[$(\Rightarrow)$] Let $(a,b)\ne\emptyset$ and choose $p\in(a,b)$, $r\defeq\frac12\min\{p-a,b-p\}$. If $p\in E$, done. If $p\notin E$, density gives $p\in E^{\operatorname{acc}}$, so choose $q\in(E\setminus\{p\})\cap B_r(p)$. Since $B_r(p)\subseteq(a,b)$, $q\in(a,b)\cap E$.
\item[$(\Leftarrow)$] Suppose $(a,b)\cap E\ne\emptyset$ for every nonempty open interval $(a,b)$.\\
For $p\in\R$, either $p\in E$, or $p\notin E$ and $\forall\varepsilon>0$, choose $q\in(p-\varepsilon,p+\varepsilon)\cap E$.\\
Then $q\ne p$ and $q\in B_\varepsilon(p)$, so $p\in E^{\operatorname{acc}}$. Hence $\R=E\cup E^{\operatorname{acc}}=\overline{E}$, so $E$ is dense in $\R$.
\end{enumerate}
\hlt{Structural Commentary}: On $\R$, balls and open intervals test the same local condition.\\
\hlt{Key Insight}: Density means no interval can avoid the set.
\end{proof}
\end{proposition}

\begin{definition} \label{def:separable-metric-space} \hlt{Separable Metric Space}\\
A metric space $X$ is \hlt{separable} if it contains a countable dense subset.
\end{definition}

\begin{example} \label{ex:separable-and-nonseparable-metric-spaces} \hlt{Separable and Nonseparable Metric Spaces}
\begin{enumerate}[label=(\roman*)]
\setlength{\itemsep}{0pt}
\item $\R$ is separable: $\Q$ is countable and dense. Hence $\R^k$ is separable via $\Q^k$, and $\C^k$ via $\C^k\simeq\R^{2k}$.
\item If $X$ is discrete, then $X$ is separable $\Longleftrightarrow X$ is countable,\\
since $B_{1/2}(x)=\{x\}$ forces every dense $D\subseteq X$ to contain $x$.
\item $\ell^\infty(\F)$ is not separable. The set $\{0,1\}^{\N}$ is uncountable, and $s\ne t$ gives $d_\infty(s,t)=1$; the balls $B_{1/3}(s)$ are pairwise disjoint, so no countable dense set can meet all of them.
\item For $1\leq p<\infty$, $\ell^p(\F)$ is separable. Let $D$ be the countable set of finite-support sequences with rational coordinates, using rational real and imaginary parts if $\F=\C$.\\
For $x\in\ell^p$ and $\varepsilon>0$, choose $N$ with $\sum_{n>N}\abs{x_n}^p<(\varepsilon/2)^p$, then choose $y\in D$ with support in $\{1,\ldots,N\}$ and $\sum_{n=1}^N\abs{x_n-y_n}^p<(\varepsilon/2)^p$. Thus $\norm{x-y}_p<\varepsilon$.
\end{enumerate}
\hlt{Structural Commentary}: Separability is the existence of a countable set meeting every nonempty open ball.\\
\hlt{Key Insight}: Finite rational data are dense in $\ell^p$, while $\ell^\infty$ contains an uncountable $1$-separated set.
\end{example}

\begin{proposition} \label{prop:compact-metric-spaces-separable} \hlt{Compact Metric Spaces Are Separable}\\
Let $(X,d)$ be a compact metric space. Then $\exists$ countable $D\subseteq X$ such that $\overline{D}^{\,X}=X$.
\begin{proof} By Theorem \ref{thm:compact-complete-totally-bounded}, $X$ is totally bounded.\\ 
$\forall n\in\N$, choose finite $F_n\subseteq X$ such that $X\subseteq\bigcup_{a\in F_n}B_{1/n}(a)$, set $D\defeq\bigcup_{n=1}^{\infty}F_n$. Then $D$ is countable.\\
Let $x\in X$ and $r>0$. Choose $n\in\N$ with $1/n<r$.\\
Since $F_n$ is a $1/n$-net for $X$, $\exists a\in F_n$ such that $d(x,a)<1/n<r$.\\
Thus $a\in B_r(x)\cap D$, so every $X$-open ball around $x$ meets $D$. Hence $x\in\overline{D}^{\,X}$, and therefore $\overline{D}^{\,X}=X$.\\
\hlt{Structural Commentary}: Compactness gives finite nets at every scale; their countable union is dense.\\
\hlt{Key Insight}: Total boundedness converts compactness into countably many finite approximations.
\end{proof}
\end{proposition}

\begin{proposition} \label{prop:metric-separable-countable-base} \hlt{Separable Metric Spaces and Countable Bases}\\
Let $(X,d)$ be a metric space, and let $\tau_d$ be its metric topology. The following are equivalent:
\begin{enumerate}[label=(\roman*)]
\setlength{\itemsep}{0pt}
\item $\exists$ countable $D\subseteq X$ such that $\overline{D}^{\,X}=X$.
\item $\exists$ countable $\mathcal{B}\subseteq\tau_d$ such that $\mathcal{B}$ is a basis for $\tau_d$.
\end{enumerate}
\begin{proof}
\begin{enumerate}[label=\roman*.]
\setlength{\itemsep}{0pt}
\item[$(i)$]$\Rightarrow (ii)$ Let $D\subseteq X$ be countable with $\overline{D}^{\,X}=X$, and set $\mathcal{B}\defeq\{B_q(x):x\in D,\ q\in\Q,\ q>0\}$.\\
Then $\mathcal{B}$ is countable. For $U\in\tau_d$ and $y\in U$, choose $r>0$ with $B_r(y)\subseteq U$, then choose $q\in\Q$ with $0<q<r/2$, $x\in D\cap B_q(y)$. Then $y\in B_q(x)\subseteq B_r(y)\subseteq U$, since $z\in B_q(x)$ implies $d(z,y)<2q<r$. QED.
\item[$(ii)$]$\Rightarrow (i)$ Let $\mathcal{B}=\{U_n:n\in I\}$ be a countable basis for $\tau_d$, with $I\subseteq\N$.\\
Choose $x_n\in U_n$ for each nonempty $U_n$, and set $D\defeq\{x_n:U_n\neq\emptyset\}$. Then $D$ is countable.\\
If $\emptyset\neq U\in\tau_d$, choose $y\in U$ and $n\in I$ with $y\in U_n\subseteq U$; then $x_n\in D\cap U$.\\
Hence every nonempty open set meets $D$, so $\overline{D}^{\,X}=X$.
\end{enumerate}
\hlt{Structural Commentary}: In metric spaces, one countable dense set produces all rational-radius local tests.\\
\hlt{Key Insight}: Separability and second countability coincide for metric spaces.
\end{proof}
\end{proposition}

\begin{proposition} \label{prop:subspaces-of-separable-metric-spaces-separable} \hlt{Subspaces of Separable Metric Spaces}\\
Let $(X,d)$ be a separable metric space and let $E\subseteq X$ with the subspace metric $d|_{E\times E}$. Then $E$ is separable.
\begin{proof}
By Proposition \ref{prop:metric-separable-countable-base}, the metric topology $\tau_d$ has a countable basis $\mathcal{B}$. Set $\mathcal{B}_E\defeq\{B\cap E:B\in\mathcal{B}\}$. Then $\mathcal{B}_E$ is countable. If $U$ is open in $E$ and $x\in U$, then $U=E\cap G$ for some $G\in\tau_d$; choose $B\in\mathcal{B}$ with $x\in B\subseteq G$. Hence $x\in B\cap E\subseteq U$, so $\mathcal{B}_E$ is a countable basis for $E$. By Proposition \ref{prop:metric-separable-countable-base}, $E$ is separable.\\
\hlt{Structural Commentary}: A countable ambient basis restricts to a countable subspace basis.\\
\hlt{Key Insight}: Metric separability passes to all subspaces through countable bases.
\end{proof}
\end{proposition}

\begin{corollary} \label{cor:euclidean-subsets-separable} \hlt{Euclidean Subsets Are Separable}\\
Let $k\in\N$ and let $E\subseteq\R^k$ have the subspace Euclidean metric. Then $E$ is separable.
\begin{proof} The set $\Q^k$ is countable and dense in $\R^k$.\\
For $x\in\R^k$ and $r>0$, choose $q=(q_1,\ldots,q_k)\in\Q^k$ with $\abs{x_i-q_i}<r/\sqrt{k}$ for each $i$, so $\norm{x-q}_2<r$.\\
Thus $\R^k$ is separable, and Proposition \ref{prop:subspaces-of-separable-metric-spaces-separable} gives that $E$ is separable.\\
\hlt{Structural Commentary}: Rational grids are countable dense nets.\\
\hlt{Key Insight}: Euclidean separability passes to subspaces.
\end{proof}
\end{corollary}

\begin{theorem} \label{thm:supremum-belongs-to-closure} \hlt{Suprema Belong to Closures}\\
Let $\emptyset\ne E\subseteq\R$ be bounded above and $y\defeq\sup E$. Then $y\in\overline{E}$; if $E$ is closed, then $y\in E$.
\begin{proof}
If $y\in E$, then $y\in\overline{E}$. If $y\notin E$, then $\forall h>0$, $y-h$ is not an upper bound; choose $x_h\in E$ with $y-h<x_h\leq y$. Thus $x_h\ne y$, $\abs{x_h-y}<h$, and $B_h(y)\cap(E\setminus\{y\})\ne\emptyset$, so $y\in E^{\operatorname{acc}}\subseteq\overline{E}$.\\
If $E$ is closed, $\overline{E}=E$ by Theorem \ref{thm:closure-smallest-closed-containing-set}(ii).\\
\hlt{Structural Commentary}: LUB gives points of $E$ arbitrarily close below $\sup E$.\\
\hlt{Key Insight}: A supremum is either attained or approached.
\end{proof}
\end{theorem}

\begin{method} \label{meth:metric-topology-proof-patterns}\label{meth:metric-topology-proof-patterns-source-todo} \hlt{Metric Proof Patterns}
\begin{enumerate}[label=(\roman*)]
\setlength{\itemsep}{0pt}
\item To prove $G$ is open, fix $x\in G$ and find $r>0$ with $B_r(x)\subseteq G$.
\item To prove $F$ is closed, show $F^{\operatorname{acc}}\subseteq F$, or equivalently show $F^c$ is open.
\item To prove $C=\overline{E}$, show $E\subseteq C$, $C$ is closed, and $C\subseteq F$ for every closed $F\supseteq E$.
\item To prove $p\in\overline{E}$, show $\operatorname{dist}(p,E)=0$, or build a sequence in $E$ converging to $p$.
\end{enumerate}
\end{method}
